% !TEX root = ../tesi.tex


% Descrive l'approccio iterativo in modo preciso e formale.
%
% 4.1 Algorithm Description
%   - Pseudocodice formale dell'approccio iterativo
%   - Input/output espliciti
%   - Aggiornamento dello stato tra un'iterazione e l'altra (initial, new, kept)
%?   - Gestione del caso UNSAT a ogni iterazione


%? 4.2 Formal Properties
%? - Da valutare se necessarario dimostrare correttezza e formalità di esecuzione; forse non necessario dato che stiamo Utilizzando Verefoo in modo "legittimo" e semplicemente lavorando su I/O di Verefoo, senza modificare il solver o la logica di base.


% 4.3 Worked Example
%   - Topologia sintetica piccola (3 nodi, 3-4 NSR)
%   - Mostra iterazione per iterazione come si muovono gli NSR tra i set initial/new/kept => far capire nel dettaglio la meccanica dell'approccio

% todo: immagine/schema della struttura dei set initial/new/kept e come cambiano tra un'iterazione e l'altra

% Sezione sull'ordinamento degli NSR (NSR Ordering): menzionare come implementazione futura

\section{Algorithm Description}

This section describes the iterative approach in detail, highlighting the structural differences with respect to vanilla VEREFOO. Both variants share the same underlying MaxSMT engine and the same formal correctness guarantees; what changes is how the set of NSRs is presented to the solver across invocations.

\subsection{Vanilla VEREFOO}

In the standard, non-incremental version of VEREFOO, all Network Security Requirements are collected into a single set and passed to the solver in one call. The algorithm proceeds as follows.

\begin{algorithm}[H]
\caption{Vanilla VEREFOO}
\label{alg:vanilla}
\begin{algorithmic}[1]
\Require Network graph $G$, full NSR set $P = \{p_1, p_2, \ldots, p_n\}$
\Ensure Firewall allocation and filtering rules satisfying all $p \in P$, or UNSAT
\State Build Allocation Graph from $G$
\State $\mathit{result} \gets \textsc{SolveMaxSMT}(G,\; \mathit{initial} = \emptyset,\; \mathit{added} = P)$
\If{$\mathit{result} = \text{SAT}$}
    \State Extract and return firewall allocation and filtering rules
\Else
    \State \Return UNSAT
\EndIf
\end{algorithmic}
\end{algorithm}

A single MaxSMT call encodes all $n$ requirements simultaneously. The solver searches for a globally optimal configuration, minimising the total number of allocated firewalls and rules across the entire NSR set.

% todo: inserire qui il codice del costruttore vanilla --- VerefooSerializer(NFV root, String algo)

\subsection{React-VEREFOO}

React-VEREFOO extends vanilla VEREFOO with the ability to reuse an existing valid configuration when the NSR set changes, rather than recomputing everything from scratch~\cite{reactverefoo}. It takes two NSR sets as input: the \textit{initial} set $P_0$, representing the requirements already enforced by the current configuration, and the \textit{added} set $R_a$, representing the new requirements that must be incorporated. Internally, React-VEREFOO classifies every NSR into deleted ($R_d$), added ($R_a$), or kept ($R_k$) groups (as described in Section~\ref{chap:background}), identifies the minimal network area affected by the added requirements, and restricts the solver to that area while leaving the rest of the configuration untouched.

\begin{algorithm}[H]
\caption{React-VEREFOO}
\label{alg:react}
\begin{algorithmic}[1]
\Require Current graph $G$, initial NSR set $P_0$, added NSR set $R_a$
\Ensure Updated firewall allocation satisfying all $p \in P_0 \cup R_a$, or UNSAT
\State Classify NSRs: $R_k \gets P_0$,\quad $R_d \gets \emptyset$,\quad $R_a \gets R_a$
\State Identify reconfigurable area of $G$ based on $R_a$
\State $\mathit{result} \gets \textsc{SolveMaxSMT}(G,\; \mathit{initial} = P_0,\; \mathit{added} = R_a)$
\If{$\mathit{result} = \text{SAT}$}
    \State Extract and return updated firewall allocation and filtering rules
\Else
    \State \Return UNSAT
\EndIf
\end{algorithmic}
\end{algorithm}

React-VEREFOO issues a single MaxSMT call, just like vanilla VEREFOO, but the problem it encodes is smaller: only the nodes in the reconfigurable area appear as free decision variables. The efficiency gain therefore depends directly on how large $R_a$ is relative to the total NSR set — the smaller the added set, the fewer nodes are flagged, and the smaller the solver problem. This relationship is captured by the percentage of requirements kept: when most requirements are kept, React-VEREFOO is significantly faster than vanilla; when most requirements change, the benefit in time execution is reduced.

React-VEREFOO thus provides the reconfiguration primitive that the iterative approach builds upon: by calling it repeatedly with $|R_a| = 1$ at each step, the iterative variant ensures that percentage of requirements kept is always at its maximum possible value, regardless of how many new NSRs must ultimately be incorporated.

\subsection{Iterative VEREFOO}

The iterative variant modifies this structure by introducing NSRs one at a time, using the React-VEREFOO reconfiguration mechanism at each step. The key insight is that at every iteration the solver receives a minimal incremental problem: one new requirement on top of an otherwise stable, already-verified configuration.

The algorithm maintains two evolving structures across iterations: $\mathit{consideredProps}$, the set of NSRs already enforced and passed as the \textit{initial} set to React-VEREFOO; and $\mathit{currentGraph}$, the network graph as updated by the most recent SAT result. Both are initialised before the loop begins: $\mathit{consideredProps}$ with any pre-existing NSRs (denoted $P_0$, which may be empty), and $\mathit{currentGraph}$ with the original input graph.

\begin{algorithm}[H]
\caption{Iterative VEREFOO}
\label{alg:iterative}
\begin{algorithmic}[1]
\Require Network graph $G$, pre-existing NSR set $P_0$, new NSR sequence $P = (p_1, p_2, \ldots, p_n)$
\Ensure Firewall allocation and filtering rules satisfying all $p \in P_0 \cup P$, or UNSAT at the failing step
\State Build Allocation Graph from $G$
\State $\mathit{consideredProps} \gets P_0$
\State $\mathit{currentGraph} \gets G$
\For{each $p_i \in P$}
    \State $\mathit{result} \gets \textsc{SolveMaxSMT}(\mathit{currentGraph},\; \mathit{initial} = \mathit{consideredProps},\; \mathit{added} = \{p_i\})$
    \If{$\mathit{result} = \text{SAT}$}
        \State $\mathit{consideredProps} \gets \mathit{consideredProps} \cup \{p_i\}$
        \State $\mathit{currentGraph} \gets$ updated graph from $\mathit{result}$
    \Else
        \State \Return UNSAT at step $i$
    \EndIf
\EndFor
\State \Return firewall allocation and filtering rules from $\mathit{currentGraph}$
\end{algorithmic}
\end{algorithm}

The total number of MaxSMT calls equals $n = |P|$, one per new NSR. Each call operates on a problem whose \textit{added} set contains exactly one requirement — the smallest possible incremental problem. The \textit{initial} set grows by one element at each step, shifting requirements from added to kept as the loop progresses and maximising PercReqKept at every iteration (as defined in Chapter~\ref{chap:objectives}).

% todo: inserire qui il codice del costruttore iterativo --- VerefooSerializer(NFV root, String algo, Boolean react)

\subsection{Structural Comparison}

The two variants differ in three structural points. First, the number of solver invocations: vanilla VEREFOO as well as the React version, issues one MaxSMT call regardless of the number of requirements, while the iterative variant issues a number of calls equal to the number of new requirements. Secondly, the scope of each call: vanilla encodes the full NSR set at once, while each iterative call encodes only one new requirement on top of a fixed base. Finally, the state passed between calls: vanilla has no inter-call state, while the iterative variant carries $\mathit{consideredProps}$ and $\mathit{currentGraph}$ across iterations to keep track of the evolving configuration, ensuring that each step preserves all previously enforced requirements.

% todo: valutare se aggiungere una tabella comparativa sintetica vanilla vs iterativo

\section{Worked Example}

To make the mechanics of the iterative approach concrete, this section traces a complete execution on a small example network. The example uses whitelisting policy mode, in which all traffic is blocked by default and only flows covered by an explicit ALLOW rule are permitted.

\subsection{Network Topology and NSR Set}

The network consists of three endpoint nodes: a \textit{Client} (C), a \textit{WebServer} (W), and a \textit{Database} (D), connected in a linear chain. Traffic from C to D must therefore transit through W. Two Allocation Places — $\mathit{AP}_1$ on the C–W link and $\mathit{AP}_2$ on the W–D link — are the candidate positions where the solver may allocate firewalls.

\begin{figure}[H]
\centering
\begin{tikzpicture}[
    node/.style={circle, draw, minimum size=1cm, font=\small},
    ap/.style={rectangle, draw, dashed, minimum size=0.8cm, font=\scriptsize},
    >=Stealth
]
\node[node] (C) {C};
\node[ap,  right=1.2cm of C] (AP1) {$\mathit{AP}_1$};
\node[node, right=1.2cm of AP1] (W) {W};
\node[ap,  right=1.2cm of W] (AP2) {$\mathit{AP}_2$};
\node[node, right=1.2cm of AP2] (D) {D};

\draw[<->] (C) -- (AP1);
\draw[<->] (AP1) -- (W);
\draw[<->] (W) -- (AP2);
\draw[<->] (AP2) -- (D);
\end{tikzpicture}
\caption{Synthetic topology used in the worked example. Dashed rectangles are Allocation Places; solid circles are endpoints.}
\label{fig:example-topology}
\end{figure}

Three NSRs must be enforced, introduced one at a time in the following order:

\begin{itemize}
  \item $p_1$: C $\to$ W — \textit{reachability} (client must be able to reach the web server)
  \item $p_2$: W $\to$ D — \textit{reachability} (web server must be able to access the database)
  \item $p_3$: C $\to$ D — \textit{isolation} (client must not be able to reach the database directly)
\end{itemize}

The initial pre-existing set is $P_0 = \emptyset$: no firewall is allocated and no requirement is enforced at the start.

\subsection{Iteration 1 — Introducing \texorpdfstring{$p_1$}{p1}}

\begin{itemize}
  \item $\mathit{initial} = \emptyset$,\quad $\mathit{added} = \{p_1\}$
  \item $R_k = \emptyset$,\quad $R_a = \{p_1\}$,\quad $R_d = \emptyset$
\end{itemize}

React-VEREFOO analyses the path C$\to$W and identifies $\mathit{AP}_1$ as the relevant allocation area. Because all traffic is blocked by default and $p_1$ requires the path C$\to$W to be permitted, the solver allocates a firewall $\mathit{FW}_1$ at $\mathit{AP}_1$ with the following filtering policy:

\begin{center}
\texttt{FW\textsubscript{1} @ AP\textsubscript{1}:\ \ ALLOW C\,$\to$\,W\ \ |\ \ default: DROP}
\end{center}

After the solver returns SAT, $p_1$ is moved into \textit{consideredProps} and the graph is updated to $G_1$.

\subsection{Iteration 2 — Introducing \texorpdfstring{$p_2$}{p2}}

\begin{itemize}
  \item $\mathit{initial} = \{p_1\}$,\quad $\mathit{added} = \{p_2\}$
  \item $R_k = \{p_1\}$,\quad $R_a = \{p_2\}$,\quad $R_d = \emptyset$
\end{itemize}

The reconfigurable area is now the path W$\to$D, which involves $\mathit{AP}_2$. The firewall $\mathit{FW}_1$ at $\mathit{AP}_1$ is kept untouched, as it lies outside the identified area. The solver allocates a second firewall $\mathit{FW}_2$ at $\mathit{AP}_2$:

\begin{center}
\texttt{FW\textsubscript{1} @ AP\textsubscript{1}:\ \ ALLOW C\,$\to$\,W\ \ |\ \ default: DROP} \\[4pt]
\texttt{FW\textsubscript{2} @ AP\textsubscript{2}:\ \ ALLOW W\,$\to$\,D\ \ |\ \ default: DROP}
\end{center}

After the execution of this iteration, $p_2$ is added to \textit{consideredProps} and the graph is updated to $G_2$.

\subsection{Iteration 3 — Introducing \texorpdfstring{$p_3$}{p3}}

\begin{itemize}
  \item $\mathit{initial} = \{p_1, p_2\}$,\quad $\mathit{added} = \{p_3\}$
  \item $R_k = \{p_1, p_2\}$,\quad $R_a = \{p_3\}$,\quad $R_d = \emptyset$
\end{itemize}

$p_3$ requires that C cannot reach D. The only path from C to D passes through $\mathit{AP}_1$. At $\mathit{AP}_1$, $\mathit{FW}_1$ has a single ALLOW rule for traffic whose destination is W; any packet from C destined for D carries a different destination address and therefore hits the default DROP action. The isolation requirement is already satisfied by the current configuration without any modification.

React-VEREFOO confirms the existing configuration as SAT and no new firewall or rule is added. $p_3$ is moved into \textit{consideredProps} and the graph remains $G_2$.

\subsection{Final Configuration and Summary}

The final allocation consists of two firewalls, one at each allocation place, with minimal rule sets:

\begin{center}
\texttt{FW\textsubscript{1} in AP\textsubscript{1}:\ \ ALLOW C\,$\to$\,W\ \ |\ \ default: DROP} \\[4pt]
\texttt{FW\textsubscript{2} in AP\textsubscript{2}:\ \ ALLOW W\,$\to$\,D\ \ |\ \ default: DROP}
\end{center}

Table~\ref{tab:worked-example} summarises the state of all sets at each iteration.

\begin{table}[H]
\centering
\caption{Evolution of NSR sets and PercReqKept across the three iterations.}
\label{tab:worked-example}
\begin{tabularx}{\textwidth}{lXXXXX}
\hline
\textbf{Iter.} & \textbf{Added} & $R_k$ & $R_a$ & $R_d$ \\
\hline
1 & $p_1$ & $\emptyset$ & $\{p_1\}$ & $\emptyset$ \\
2 & $p_2$ & $\{p_1\}$ & $\{p_2\}$ & $\emptyset$ \\
3 & $p_3$ & $\{p_1, p_2\}$ & $\{p_3\}$ & $\emptyset$ \\
\hline
\end{tabularx}
\end{table}

Three observations follow from this trace. First, as more requirements are incorporated, the fraction of the network that the solver must reconsider is reduced for each iteration: at each step, the nodes satisfying all previously introduced requirements are treated as fixed, and only the area relevant to the single new requirement is left as a free decision variable. Secondly, the already-enforced requirements impose no additional solving cost — their satisfaction is guaranteed by construction from the previous step and the solver does not need to reason about them again. Third, iteration 3 required no modification at all: the isolation requirement was already implied by the whitelisting default of $\mathit{FW}_1$, demonstrating that the iterative approach can incorporate new NSRs with zero reconfiguration cost when the existing configuration already covers them.

It should be noted that these results are a consequence of the deliberately small scale of this example. In a real network, with a larger number of nodes, more complex topologies, and a higher number of NSRs, the interactions between requirements are significantly more complex and the outcomes of each iteration less predictable. Three dedicated testing scenarios on real and synthetic topologies are reported and discussed in Chapter~\ref{chap:implementation}.
